For each source and order write
\[
 \lambda_{j,\nu}=\operatorname{cum}_\nu(\varepsilon_j).
\]
The moment-determinacy assumption entails that the moment sequence, and hence every cumulant used below, exists. To derive the observable tensor formula without adding an independence property silently, recall the moment--cumulant identity
\[
 \operatorname{cum}(Z_1,\ldots,Z_\nu)
 =\sum_{\pi\in\mathfrak P_\nu}
 (|\pi|-1)!(-1)^{|\pi|-1}
 \prod_{B\in\pi}\mathbb E\!\left[\prod_{s\in B}Z_s\right].
\]
This finite polynomial is multilinear. If its arguments split into two nonempty independent subfamilies, factorization of the moments in the display and grouping the partitions by their induced partitions on the two subfamilies makes the alternating partition sum zero. Therefore, after inserting \(X_i=\sum_j C_{ij}\varepsilon_j\), mutual independence gives
\[
 \mathcal T_\nu:=\operatorname{cum}_\nu(X)
 =\sum_{j=1}^n\lambda_{j,\nu}c_j^{\otimes\nu}
 =\sum_{j\in A_\nu}\lambda_{j,\nu}c_j^{\otimes\nu}.
 \tag{1}
\]

Fix \(\nu\ge2n-1\) and set \(s=n-1\) and \(t=\nu-2n+2\), so \(s\ge1\) because \(n\ge p\ge2\), and \(t\ge1\). Reshaping (1) into blocks of degrees \((s,s,t)\) yields
\[
 \widetilde{\mathcal T}_\nu
 =\sum_{j\in A_\nu}\lambda_{j,\nu}
 a_j\otimes b_j\otimes d_j,
 \quad
 a_j=b_j=\operatorname{vec}(c_j^{\otimes s}),
 \quad
 d_j=\operatorname{vec}(c_j^{\otimes t}).
 \tag{2}
\]
Let \(R=R_\nu\). By pairwise nonproportionality of the columns and \(s=n-1\ge R-1\), the already established Veronese--Lagrange rank lemma gives Kruskal rank \(R\) to each of the first two factor matrices in (2). If \(R\ge2\), every pair of third-mode columns is linearly independent. Indeed, if \(z^{\otimes t}\) and \(w^{\otimes t}\) were proportional, then for every linear functional \(f\) with \(f(w)=0\), contraction by \(f^{\otimes t}\) would give \(f(z)^t=0\); hence every functional vanishing on \(w\) would vanish on \(z\), forcing \(z\in\operatorname{span}(w)\). Thus the third factor matrix has Kruskal rank at least two. Consequently
\[
 k_1+k_2+k_3\ge R+R+2=2R+2.
\]
The already established three-factor uniqueness lemma therefore shows that (2) is unique up to a common permutation and reciprocal nonzero factor scalings.

The number \(R\) in (2) is also the minimal tensor rank, rather than merely the length of one decomposition. The first factor matrix \(A=[a_j]_{j\in A_\nu}\) and the second factor matrix \(B=[b_j]_{j\in A_\nu}\) have full column rank. The Khatri--Rao matrix \([b_j\otimes d_j]_{j\in A_\nu}\) has full column rank: if \(\sum_j q_j b_j\otimes d_j=0\), applying the \(j\)-th row of a left inverse of \(B\) to the first tensor component gives \(q_jd_j=0\), and \(d_j\ne0\) gives \(q_j=0\). Hence the first matricization of (2),
\[
 A\operatorname{diag}(\lambda_{j,\nu})
 [b_j\otimes d_j]_{j\in A_\nu}^{\top},
\]
has matrix rank \(R\). A sum of fewer than \(R\) rank-one tensors would have a first matricization of rank less than \(R\), which is impossible. Thus the constrained minimization in the corrected multi-order handle has value \(R_\nu^\star=R_\nu\). Its feasible decomposition (2) is a minimal unrestricted decomposition, so uniqueness shows that every minimizer returns precisely \(\{[c_j]:j\in A_\nu\}\). The cases \(R=0\) and \(R=1\) are direct: the tensor is respectively zero or a nonzero rank-one tensor, and in the latter case the first factor has the unique projective Veronese preimage \([c_j]\). The contraction argument above also proves that the projective Veronese map is injective, so extracting that preimage is unambiguous.

The same calculation yields the advertised adaptive certificate. If \(R_\nu\ge2\), \(\nu\ge2R_\nu-1\), and the block degrees are
\[
 (R_\nu-1,R_\nu-1,\nu-2R_\nu+2),
\]
then the first two degrees are at least \(R_\nu-1\) and the third is at least one. Their Kruskal ranks are therefore \(R_\nu,R_\nu,\) and at least two, respectively. The same Kruskal inequality and matricization-rank argument prove minimality, uniqueness, and exact recovery. This establishes the adaptive threshold without requiring it for the universal construction, whose blocks depend only on the known \(n\).

By the already established Marcinkiewicz cumulant-tail lemma, each nondegenerate non-Gaussian source having moments of all orders has infinitely many nonzero cumulants of order at least three. In particular, for every \(j\) there exists a finite \(\nu_j\ge2n-1\) with \(\lambda_{j,\nu_j}\ne0\). Therefore
\[
 K_{C,\varepsilon}:=\max_{1\le j\le n}\nu_j<\infty.
 \tag{3}
\]
Processing successive orders and taking the set-theoretic union of their recovered projective directions reaches cardinality \(n\) no later than (3). Pairwise nonproportionality makes exact projective deduplication valid, and no label is transported between orders.

Choose any nonzero representative of each recovered direction and form a mixing matrix \(\widetilde C\). It differs from \(C\) by a column permutation and nonsingular column scaling. Such scaling and permutation preserve whether a treatment loading is zero, preserve every deletion rank because right multiplication of \(C_{-x,-j}\) by an invertible diagonal matrix does not change rank, and preserve each ratio \(C_{ij}/C_{xj}\). Hence the observable recovery map returns \(\mathcal V_x(C)\) and \(\mathcal R_{xy}(C)\). These are not being defined as causal effects: the already established sharp effect-frontier theorem proves that they equal, respectively, the law-fiber response-vector target and its \(y\)-coordinate projection.

We now prove that no cumulant cutoff can be uniform on the qualitative class. Fix an arbitrary finite \(K\ge3\). The already established finite-moment near-Gaussian witness lemma supplies, for any \(\rho>0\), a centered, variance-one, non-Gaussian, moment-determinate law \(F\) whose moments through order \(K\) equal those of a standard Gaussian. Let the two sources be independent with common law \(F\), and fix \(0<\gamma<\pi/4\). Consider
\[
 C_0=I_2,
 \qquad
 C_1=
 \begin{pmatrix}
 \cos\gamma&-\sin\gamma\\
 \sin\gamma&\cos\gamma
 \end{pmatrix}.
 \tag{4}
\]
Both matrices have full row rank and normalized, nonzero, pairwise nonproportional columns. Independence makes every joint source moment of total degree at most \(K\) equal to the corresponding moment of \(\mathcal N(0,I_2)\). Since \(C_1\) is orthogonal and the Gaussian law is orthogonally invariant, the two observed models in (4) have identical moments, and therefore identical cumulant tensors, through order \(K\). Their projective direction sets are distinct. Thus no rule receiving cumulants only through a fixed order \(K\) can recover both, proving that the stopping cutoff in (3) is necessarily representation dependent.

The same pair gives the statistical impossibility. For each \(N\), choose the witness law \(F_N\) so that
\[
 d_{\mathrm{TV}}(F_N,\mathcal N(0,1))\le\frac1{4N^2},
\]
and let \(P_{k,N}\) be the observed law generated by \(C_k\) and two independent \(F_N\) sources. Total variation contracts under a measurable pushforward directly from its supremum-over-events definition. Also, telescoping signed product measures gives
\[
 d_{\mathrm{TV}}(\mu^{\otimes m},\nu^{\otimes m})
 \le m\,d_{\mathrm{TV}}(\mu,\nu).
 \tag{5}
\]
Applying these facts twice and inserting the common orthogonally invariant law \(\mathcal N(0,I_2)\) gives
\[
 d_{\mathrm{TV}}(P_{0,N},P_{1,N})
 \le4d_{\mathrm{TV}}(F_N,\mathcal N(0,1))
 \le N^{-2},
\]
and then (5) yields
\[
 d_{\mathrm{TV}}(P_{0,N}^{\otimes N},P_{1,N}^{\otimes N})
 \le N^{-1}\longrightarrow0.
 \tag{6}
\]

The two finite projective direction sets in (4) have a fixed positive Hausdorff separation. A uniformly consistent direction-set estimator would, by choosing the closer of these two targets, induce a test with both error probabilities tending to zero. Yet for every test region \(E\), the definition of total variation gives
\[
 P_{0,N}^{\otimes N}(E^c)+P_{1,N}^{\otimes N}(E)
 \ge1-d_{\mathrm{TV}}(P_{0,N}^{\otimes N},P_{1,N}^{\otimes N}),
\]
which tends to one by (6), a contradiction.

The causal targets are separated as well. For \(x=1\) and \(y=2\), direct evaluation of the one-dimensional deletion matrices gives
\[
 \mathcal R_{12}(C_0)=\{0\},
 \qquad
 \mathcal R_{12}(C_1)=\{\tan\gamma,-\cot\gamma\}.
 \tag{7}
\]
Their Hausdorff distance is positive, so the same nearest-target test rules out uniform consistency for the frontier. More strongly, let \(\mathcal U_N\) be any random compact subset of \(\mathbb R\) satisfying uniform frontier coverage at least \(1-\beta\), where \(\beta<1/2\). Under the first experiment the event \(\mathcal R_{12}(C_0)\subseteq\mathcal U_N\) has probability at least \(1-\beta\). By coverage under the second experiment and (6), under the first experiment the event \(\mathcal R_{12}(C_1)\subseteq\mathcal U_N\) has probability at least \(1-\beta-o(1)\). Their intersection therefore has probability at least \(1-2\beta-o(1)>0\). On that intersection,
\[
 d_H(\mathcal U_N,\mathcal R_{12}(C_0))
 \ge \sup_{z\in\mathcal R_{12}(C_1)}|z|
 =\cot\gamma.
\]
Thus no uniformly honest confidence union of level strictly greater than one half has uniformly vanishing Hausdorff error.

Finally, the causal map has two further deterministic discontinuity boundaries. Column normalization does not alter ratios or deletion ranks, so normalize the following columns if desired. For
\[
 D_t=\begin{pmatrix}t&1\\1&1\end{pmatrix},
\]
the matrices are full rank with nonzero nonproportional columns near \(t=0\), and
\[
 \mathcal R_{12}(D_t)=\{1/t,1\}\quad(t\ne0),
 \qquad
 \mathcal R_{12}(D_0)=\{1\}.
 \tag{8}
\]
This is the vanishing-denominator boundary. For
\[
 E_t=
 \begin{pmatrix}
 1&1&0&1\\
 1&t&0&0\\
 0&0&1&1
 \end{pmatrix},
\]
the same rank and direction conditions hold near zero, while explicit deletion gives
\[
 \mathcal R_{12}(E_0)=\{0\},
 \qquad
 \mathcal R_{12}(E_t)=\{0,t,1\}\quad(t\ne0).
 \tag{9}
\]
Indeed, at \(t=0\) only columns two and four are admissible and both ratios are zero; for \(t\ne0\), columns one, two, and four are admissible and give \(1,t,0\). This is the deletion-rank boundary.

Equations (6), (8), and (9) prove that qualitative moment determinacy and qualitative nonvanishing/rank conditions alone supply neither a uniform statistical modulus nor continuity of the causal map. Quantitative non-Gaussianity, moment, Veronese-conditioning, angular, probe/eigen-gap, denominator, and deletion-rank controls are therefore not consequences of the qualitative model and cannot simply be omitted from a uniform root-\(N\) claim. On the separated common-order subclass these controls are expressly assumed, and the already established common-order confidence theorem supplies simultaneous response-vector and scalar coverage with \(O(N^{-1/2})\) Hausdorff error. This completes both the positive pointwise population construction and the negative uniform boundary.